%%
%% fall-lecture11.tex
%% 
%% Made by Alex Nelson <pqnelson@gmail.com>
%% Login   <alex@lisp>
%% 
%% Started on  2026-10-22T12:17:17-0700
%% Last update 2026-10-22T12:17:17-0700
%% 

\lecture[Banach spaces]{}

\begin{definition}
Let $X$ be a vector space over $\RR$ or $\CC$. A \define{Norm} on $X$
is a map $\|-\|\colon X\to\RR$ such that
\begin{enumerate}
\item \textsc{Triangle inequality}: for all $x,y\in X$, $\|x+y\|\leq\|x\|+\|y\|$
\item \textsc{Homogenous}: for all $\lambda\in\FF$ and for all $x\in X$,
  $\|\lambda x\|=|\lambda|\cdot\|x\|$
\item \textsc{Positive-definite}: for all $x\in X$, $\|x\|\geq0$ and
  $\|x\|=0\iff x=0$.
\end{enumerate}
We call the pair $(X,\|-\|)$ a \define{Normed Vector Space}.
\end{definition}

\begin{definition}
Let $(X,\|-\|)$ be a normed vector space. We can construct the function
\begin{align}
\rho\colon & X\times X\to\RR\\
& (x,y)\mapsto\rho(x,y)=\|x-y\|\nonumber
\end{align}
which is a metric on $X$. It induces the \define{Norm Topology} on $X$.
Without loss of generality, we consider $(X,\|-\|)$ as a metric space
with this metric $\rho$ and equipped with the norm topology.
\end{definition}

\begin{definition}
We say $(X,\|-\|)$ a \define{Banach Space} if it is a complete metric
space under $\rho$.
\end{definition}

\begin{remark}
What do we do with an \emph{incomplete} normed space $(X,\|-\|)$?
Well, recall, we can always embed a metric space into a complete space
(by taking its completion). We can do likewise here, and embed any
normed space into a Banach space.
\end{remark}

\begin{definition}
Let $(X,\|-\|)$ be a normed space, let $(x_{n})$ be a sequence on $X$.
A series $\sum^{\infty}_{n}x_{n}$ is called \define{Absolutely Continuous}
if $\sum\|x_{n}\|<\infty$.
\end{definition}

\begin{theorem}
A normed space $(X,\|-\|)$ is complete if and only if every absolutely
continuous series converges.
\end{theorem}

\begin{proof}
\forwardproof\ Assume $(X,\|-\|)$ is complete. Let $\sum x_{n}$
be an arbitrary series on $X$. Assume $\sum\|x_{n}\|<\infty$. Consider
the partial sums
\begin{equation}
S_{n}=\sum^{n}_{k=1}x_{k}.
\end{equation}
Then $S_{n}$ forms a Cauchy sequence since
\begin{equation}
\|S_{n}-S_{m}\|\leq\sum^{n}_{k=m+1}\|x_{k}\|\xrightarrow{m,n\to\infty}0
\end{equation}
by absolute continuity. Then, by completeness of $X$, $S_{n}$ must
converge. Hence $\sum x_{n}$ must converge.

\backwardproof\ Assume every absolutely continuous series
converges. Let $(x_{n})$ be a Cauchy sequence in $X$. [We want to show
$(x_{n})$ converges to some point in $X$.]
Choose numbers $n_{1}<n_{2}<\dots$ such that for all $n>n_{j}$ and
$m>n_{j}$ we have $\|x_{n}-x_{m}\|\leq2^{-j}$. Set $y_{1}=x_{n_{1}}$,
$y_{2}=x_{n_{2}}-x_{n_{1}}$, \dots, $y_{k+1}=x_{n_{k+1}}-x_{n_{k}}$,
\dots. Then
\begin{equation}
\sum^{k}_{j=1}y_{j}=x_{n_{k}}
\end{equation}
because this is a telescoping series. We see $\sum_{j}^{\infty}\|y_{j}\|<\infty$ is finite
since
\begin{subequations}
  \begin{align}
\sum_{j}\|y_{j}\| &= \|x_{n_{1}}\| + \sum_{j=2}\|x_{n_{j+1}}-x_{n_{j}}\|\\
&\leq\sum_{j}2^{-j}=1.
  \end{align}
\end{subequations}
Then this estimates $\sum_{j}\|y_{j}\|<\infty$ which means the partial
sums $\sum y_{j}$ converges. This means $\lim_{n_{k}\to\infty} x_{n_{k}}$ exists
and since $(x_{n})$ is Cauchy it follows that $x_{n}$ has a
limit. Hence every Cauchy sequence converges.
\end{proof}

\begin{example}[Banach spaces]
If $K$ is a compact set, then $C(K)$ the space of continuous
real-valued functions is a vector space over $\RR$ and we can define
$\|f\|=\max_{K}|f|$. This is a complete spacde, hence Banach. (This is
especially interesting when $K$ is also Hausdorff.)
\end{example}

\begin{example}
The space of sequences whose limit at infinity is zero $(C_{0})$ is a
subspace of sequences where there exists a limit at infinity $(C)$. We
can identify $(C)$ with $C(\widehat{\NN})$ the continuous functions on
the one-point compactification of $\NN$ where $(x_{n})$ converges
$x_{n}\to x_{\infty}$ to the value assigned at infinity. This can be
viewed as a special case of the previous example since $\widehat{\NN}$
is compact. Here $(C_{0})$ is a closed subset of $(C)$.
\end{example}

\begin{example}
The space $\ell^{1}(\NN)$ consists of absolutely convergent sequences
\begin{equation}
\ell^{1}(\NN)=\{(x_{n})\mid\sum|x_{n}|<\infty\}.
\end{equation}
The norm here is
\begin{equation}
\|(x_{n})\|=\sum_{n}|x_{n}|.
\end{equation}
The claim is that it is a Banach space.
\end{example}

\begin{example}
The space 
\begin{equation}
\ell^{2}=\{(x_{n})\mid \sum_{n}|x_{n}|^{2}<\infty\}.
\end{equation}
The norm
\begin{equation}
\|(x_{n})\|=\left(\sum_{n}|x_{n}|^{2}\right)^{1/2}.
\end{equation}
It is similar to the Euclidean norm.
\end{example}

\begin{example}
More generally, 
\begin{equation}
\ell^{\infty}=\{(x_{n})\mid\max_{n}|x_{n}|<\infty\}.
\end{equation}
Its norm is
\begin{equation}
\|(x_{n})\|=\sup_{n}|x_{n}|.
\end{equation}
\end{example}

\begin{non-example}
The space $C^{1}[a,b]$ of continuously differentiable functions on the
closed interval $[a,b]$. It is a subspace of $C[a,b]$ (the space of
continuous real-valued functions on the interval), but it is not
closed under norm. Although $C[a,b]$ \emph{is} complete, the subspace
$C^{1}[a,b]$ \emph{is not} complete. (What is the completion of $C^{1}[a,b]$?
It is precisely $C[a,b]$ since $C^{1}[a,b]$ is a dense subset.)
\end{non-example}

\subsection{Linear maps}

\begin{definition}
Let $(X,\|-\|_{X})$ and $(Y,\|-\|_{Y})$ be normed spaces. Let $T\colon X\to Y$
be a linear map. We say $T$ is \define{Bounded} if there exists a
$C>0$ such that for all $x\in X$ we have $\|Tx\|_{Y}\leq C\|x\|_{X}$
(``it takes bounded sets to bounded sets'').
\end{definition}

\begin{notation}
We denote by $\mathcal{L}(X,Y)$ the set of all bounded linear maps
from $X$ to $Y$.
\end{notation}

\begin{example}
The differentiation operator $f\mapsto f'$ is an unbounded operator.
\end{example}

\begin{theorem}
Let $T\colon X\to Y$ be a linear operator. The following are equivalent:
\begin{enumerate}
\item $T$ is continuous
\item $T$ is continuous at 0
\item $T$ is bounded.
\end{enumerate}
\end{theorem}

\begin{theorem}
Let $X$ and $Y$ be normed vector spaces. Then $\mathcal{L}(X,Y)$ is a
vector space with norm
\begin{subequations}
  \begin{align}
\|T\|_{\text{op}} &= \inf\{C\in[0,+\infty)\mid\forall x\in X\ldotp\|Tx\|_{Y}\leq C\|x\|_{X}\}\\
&= \sup_{x\neq0}\frac{\|Tx\|_{Y}}{\|x\|_{X}}\\
&= \sup_{\|x\|=1}\|Tx\|_{Y}.
  \end{align}
\end{subequations}
This is called the \define{Operator Norm}. (We have the useful
estimate $\|Tx\|_{Y}\leq\|T\|_{\text{op}}\|x\|_{X}$.)
\end{theorem}

\begin{theorem}
If $Y$ is a Banach space, then $\mathcal{L}(X,Y)$ is a Banach space.
\end{theorem}

\begin{proof}
Assume $Y$ is complete. Let $(T_{n})$ be a Cauchy sequence of linear
operators from $X$ to $Y$ (i.e., in $\mathcal{L}(X,Y)$). For every
$x\in X$,
the sequence $(T_{n}x)$ is a Cauchy sequence in $Y$. (Really? Yes,
observe $\|T_{n}x\|\leq\|T_{n}\|_{\text{op}}\|x\|_{X}$ which implies $\|T_{n}x-T_{m}x\|\leq\|T_{n}-T_{m}\|_{\text{op}}\|x\|_{X}$.)
Since $Y$ is complete, $T_{n}x\to Tx$ pointwise defines the function $T$.%$T\in\mathcal{L}(X,Y)$.
Then the claim is that $T_{n}\to T$ in operator norm (this is called
\define{Strong Operator Convergence} that $\|T_{n}-T\|_{\text{op}}\to0$).
There are several steps we need to prove:
\begin{enumerate}
\item $T\colon X\to Y$ is a linear map
\item $T$ is bounded using operator norm
\end{enumerate}
\textsc{Claim}: $T_{n}\to T$ in operator norm. Let $\varepsilon>0$,
then there exists $N\in\NN$ such that for all $m,n\in\NN$ such that
$m\geq N$ and $n\geq N$ we have $\|T_{n}-T_{m}\|_{\text{op}}<\varepsilon$
by $(T_{n})$ being Cauchy. Then for every $x\in X$ such that
$\|x\|_{X}\leq1$ we have $\|T_{n}x-Tx\|\leq\varepsilon$ for all $n\geq N$.
Then take the supremum over $x\in X$ in unit ball $\|x\|_{X}\leq1$ and
it gives us $\|T_{n}-T\|_{\text{op}}<\varepsilon$.
\end{proof}

\begin{theorem}
Let $X$, $Y$, $Z$ be normed spaces.
Let $T\colon X\to Y$ and $S\colon Y\to Z$ be bounded linear maps.
Then $\|S\circ T\|_{\text{op}}\leq\|S\|_{\text{op}}\|T\|_{\text{op}}$.
In particular, if $X$ is a Banach space, then $\mathcal{L}(X,X)$ is a
Banach algebra.
\end{theorem}

(Next time, the three fundamental principles of functional analysis:
the Hahn--Banach theorem [about continuous extensions of linear
  forms], the Open Mapping Theorem, and the Banach--Steinhaus Theorem [a.k.a., Uniform Boundedness Principle].)